% ===================================================================
% CHAPTER 6: FIRST-ORDER PRIMITIVES -- CRITICAL POINTS
% ===================================================================
\chapter{First-Order Primitives: Critical Points}
\label{ch:first_order}

\section{Introduction and Related Work}
\label{sec:ch6:intro}

% New text: chapter framing against Chapter 5.
This chapter is the first rung above the zeroth-order primitives of
Chapter~\ref{ch:zeroth}. At zeroth order, a robot team reduces its
measurements to a single scalar, a mean heading or a flow magnitude, and
steers by that reduction alone. That policy reports a direction, not a
location. It carries no representation of where the field vanishes, so it
cannot certify arrival at a feature, and when applied to station-keeping
around a rotational feature it drifts outward under uncorrected centrifugal
error (Section~\ref{sec:ch6:comparison}). This chapter fits an affine model
to the positions and field measurements of three robots, one Jacobian per
time step, and recovers both the location of the nearby critical point and,
from the same fit, its topological type. The resulting estimate
$\mathbf{p}^*$ supports two control primitives: a proportional law that
drives the cluster onto $\mathbf{p}^*$, and an orbital law that holds a
commanded standoff radius around it.

% Ported from: cwaight_systems_paper.tex Sec. I (contributions list, lines
% 65-67), adapted to chapter framing.
This chapter's contributions are four. It presents a distributed algorithm
that estimates critical point location from three robots' instantaneous
measurements, with no prior map of the field. It proves that three robots
are the minimum number for which this estimation is well posed. It verifies
in simulation that the resulting estimate is sufficient to drive attraction
and orbital navigation across six canonical vector field types. It
validates both primitives in 157 hardware attraction trials and 12 hardware
orbital trials on the testbed of Chapter~\ref{ch:tools}.

% Ported from: cwaight_systems_paper.tex Sec. I, paragraphs 3-6 (lines 53,
% 57-61), trimmed to critical-point-specific material; shared adaptive
% navigation and artificial-vector-field literature moved to Chapter 1
% (\ref{ch:intro}).
By contrast, only a few works address navigation of physically sensed
vector fields \cite{55,39,34}. Work that senses the physical
field falls into two camps. One camp presumes prior knowledge of the
environment, unavailable in the unmapped settings that motivate this
chapter \cite{32,33}. The other operates online in two behaviors:
minimally actuated drifters that ride and station-keep along ambient
streamlines, and PIM-triple-inspired teams that straddle a saddle and track
a single stable or unstable manifold from local velocity measurements
\cite{34,1,37}. Both supply a heading along the flow or track
the manifold of a single assumed saddle; neither recovers the coordinate of
a critical point together with its identity across the full taxonomy of
feature types. Estimating the locations of critical points in vector
fields remains an open problem for multirobot systems.

Recovering those coordinates is harder than following a heading, for
reasons of both information and implementation. The information problem is
that the location and classification of a critical point are absent from
any single velocity sample; they are carried by the spatial derivatives of
the field \cite{56}. A single platform must move, sample again, and
difference its readings, which conflates the spatial gradient with the
field's change in time and can return features that are not physically
present. A distributed team removes this temporal aliasing by sampling at
one instant \cite{39,34,42}, but recovering the derivatives is
still not recovering the feature. Magnitude-based methods \cite{38}
reduce the flow to a scalar and yield a descent direction, but the
magnitude vanishes at every critical point regardless of type, and
estimating the two velocity components as independent scalar fields still
does not locate the point. Both the coordinate and the identity are
properties of the assembled Jacobian: inverting it recovers the location
where the field vanishes, and its eigenvalues determine the type
\cite{2}.

The implementation problem is that the local field must be reconstructed
on physical hardware from noisy distributed measurements. The conditioning
of this reconstruction degrades as the formation approaches collinearity,
so the sampling shape must be held by a formation controller under real
actuator dynamics and communication latency (Chapter~\ref{ch:estimation}).
Both the estimator and the mechatronics of the cluster limit the accuracy
of the recovered features.

\section{Critical Point Estimation from Three Robots}
\label{sec:ch6:estimation}

\subsection{Formation Model and Closed-Form Estimate}
\label{sec:ch6:formation_model}

% Ported from: cwaight_systems_paper.tex Sec. II-B, eqs. 2-9 (lines
% 107-142); thesis.tex "Distributed Critical Point Estimation"
% (lines 1909-1951).
Chapter~\ref{ch:estimation} states the local-polynomial field estimator,
its formation matrix, and its minimality condition for a general fit
order. This section specializes that machinery to fit order one, the
affine case, fit from three robots.

Given three robots sampling a planar vector field with positions
$\mathbf{p}_i = (x_i, y_i)$ and field measurements $(u_i, v_i)$ for
$i \in \{1,2,3\}$, the two field components are approximated over the
cluster's spatial extent by linear planar models
\begin{align}
U(x,y) &= ax + by + c, \label{eq:ch6:affine_u}\\
V(x,y) &= dx + ey + f. \label{eq:ch6:affine_v}
\end{align}
The coefficients solve
\begin{equation}
\mathbf{A}\boldsymbol{\theta}_u = \mathbf{u}, \qquad
\mathbf{A}\boldsymbol{\theta}_v = \mathbf{v},
\label{eq:ch6:coeff_system}
\end{equation}
where
\begin{equation}
\mathbf{A} = \begin{bmatrix} x_1 & y_1 & 1 \\ x_2 & y_2 & 1 \\ x_3 & y_3 & 1 \end{bmatrix},
\quad
\boldsymbol{\theta}_u = \begin{bmatrix} a \\ b \\ c \end{bmatrix},
\quad
\boldsymbol{\theta}_v = \begin{bmatrix} d \\ e \\ f \end{bmatrix},
\label{eq:ch6:formation_matrix}
\end{equation}
and $\mathbf{u} = [u_1,u_2,u_3]^{\mathsf T}$,
$\mathbf{v} = [v_1,v_2,v_3]^{\mathsf T}$ are the measurement vectors.

The critical point $\mathbf{p}^* = (x^*,y^*)$ satisfies $U(x^*,y^*) = 0$
and $V(x^*,y^*)=0$, which in matrix form reads
$\mathbf{J}\mathbf{p}^* = -\mathbf{h}$ with
\begin{equation}
\mathbf{J} = \begin{bmatrix} a & b \\ d & e \end{bmatrix}, \qquad
\mathbf{h} = \begin{bmatrix} c \\ f \end{bmatrix},
\label{eq:ch6:jacobian_offset}
\end{equation}
so that, provided $\det(\mathbf{J}) \neq 0$,
\begin{equation}
\mathbf{p}^* = -\mathbf{J}^{-1}\mathbf{h}.
\label{eq:ch6:pstar}
\end{equation}
The eigenvalues of $\mathbf{J}$ then identify the type of the nearby
critical point (Chapter~\ref{ch:preliminaries}).

If the true field is affine over the cluster's spatial extent,
\eqref{eq:ch6:affine_u}--\eqref{eq:ch6:affine_v} reproduce it without
error and the recovery \eqref{eq:ch6:pstar} is exact at any formation
size. On a smooth nonlinear field, the fit carries a first-order
linearization error that scales with the local curvature of the field;
this error decreases as the formation contracts toward the critical
point.

\subsection{Three Robots Are Minimal and Sufficient}
\label{sec:ch6:minimality}

% Ported from: cwaight_systems_paper.tex Sec. II-C (lines 145-153);
% thesis.tex "Minimum Robot Requirements" (lines 1952-1969).
This is the fit-order-one instance of the general minimality statement of
Chapter~\ref{ch:estimation}. Equations~\eqref{eq:ch6:affine_u}--%
\eqref{eq:ch6:affine_v} each carry three unknown coefficients, and each
robot contributes one equation to each system: robot $i$ supplies
$u_i = ax_i+by_i+c$ for the $U$-component and $v_i = dx_i+ey_i+f$ for the
$V$-component. Determining three unknowns uniquely requires at least
three equations, so three robots are the minimum for critical point
estimation. With fewer than three robots, $\mathbf{A}$ in
\eqref{eq:ch6:coeff_system} has fewer rows than columns, admitting
infinitely many coefficient solutions and preventing unique
identification of $\mathbf{p}^*$. With more than three robots the system
is overdetermined and solved by least squares.

Three non-collinear robots are also sufficient: $\det(\mathbf{A})$ equals
twice the signed area of the formation triangle, so $\mathbf{A}$ is
invertible exactly when the robots are not collinear. This is the
fit-order-one case of the common-curve degeneracy of
Chapter~\ref{ch:estimation}: a formation degenerates when its robots lie
on a common curve of the degree the fit requires, a line at first order.

\section{Critical Point Classification}
\label{sec:ch6:classification}

% Ported from: cwaight_systems_paper.tex Sec. II-A (lines 77-79);
% thesis.tex "Vector Field Environments and Critical Points"
% (lines 1863-1879, Hartman-Grobman caveat). New paragraph: planned
% hardware classification-accuracy analysis.
Chapter~\ref{ch:preliminaries} classifies a critical point by the
eigenvalue structure of the field Jacobian $\mathbf{J}$ evaluated there.
The three-robot estimator of Section~\ref{sec:ch6:formation_model}
recovers $\mathbf{J}$ directly as part of the same fit that locates
$\mathbf{p}^*$, so a type classification is available at the same
measurement cost as the location estimate, with no additional sampling.

This classification holds rigorously for hyperbolic critical points. The
center is the exception: a nonlinear perturbation of a linear center can
produce a spiral, since the Hartman-Grobman theorem does not apply at
$\alpha_{\text{eig}}=0$, and under measurement noise the estimated real
part of a center's eigenvalue hovers near zero, making
center-versus-weak-spiral discrimination ill-posed from noisy data.
Eigenvalue analysis also identifies the degenerate case
$\det(\mathbf{J})=0$, corresponding to a non-hyperbolic critical point or
a region of near-uniform flow; this case is detectable through the same
estimate but is not addressed further in this chapter.

% New analysis, 2026-09-29: experiments/classify_from_hardware_logs.py,
% results in experiments/outputs/classify_hw/summary.json (VF_Robot).
% Numbers below are from that run and await the author's review.
\subsection{Classification on Hardware}
\label{sec:ch6:classification_hw}

The hardware trials of Section~\ref{sec:ch6:attraction_hw} were designed
to test location, but every control cycle also produced a full Jacobian
estimate. Recomputing that estimate offline from the logged robot
positions and sensed hue and saturation reproduces the controller's
logged critical point estimates to within $5\times10^{-7}$~m, so the
Jacobians analyzed here are the ones the controller computed. The logs
hold 78 saddle trials and 80 vortex trials of 101 cycles each (10~s at
10~Hz).
%% TODO(reconcile): the logs contain 80 vortex runs; the critical-points
%% paper reports 79 (157 total), and its bias and precision statistics,
%% quoted elsewhere in this chapter, use that count.

Each estimate was classified by the rules of Chapter~\ref{ch:preliminaries}:
a saddle when $\det\hat{\mathbf{J}} < 0$, a node when the eigenvalues are
real and of one sign, and a rotational type when they form a complex pair
$\alpha_{\text{eig}} \pm i\omega$. Because a measured center never has
$\alpha_{\text{eig}}$ exactly zero, a complex pair was labeled a center
when $|\alpha_{\text{eig}}| < 0.2\,\omega$ and a spiral otherwise; the
threshold is a choice, and results are given for $0.1$ and $0.3$ as well.

\begin{table}[h]
\centering
\caption{Offline classification of the hardware Jacobian estimates. The
terminal window is the final 2~s of each trial; the trial vote is the
majority label over that window.}
\label{tab:ch6:classification_hw}
\begin{tabular}{lcc}
\toprule
 & Saddle map & Vortex map \\
\midrule
Trials & 78 & 80 \\
Class correct (saddle, rotational, node), all cycles & 99.9\% & 99.9\% \\
Class correct, terminal window & 100\% & 100\% \\
Class correct, trial vote & 78/78 & 80/80 \\
Exact type, terminal window, threshold 0.1 & 100\% & 71\% \\
Exact type, terminal window, threshold 0.2 & 100\% & 89\% \\
Exact type, terminal window, threshold 0.3 & 100\% & 99\% \\
\bottomrule
\end{tabular}
\end{table}

The estimate names the class of both fields correctly in effectively
every cycle, including the approach from more than 0.8~m away, so the
affine model's truncation error does not reach the classification in
these fields. Every vortex error is a stable spiral with a small negative
real part; none is a saddle or a node. This is the center-versus-spiral
ambiguity described above, and the terminal window resolves it for most
trials once the threshold exceeds the residual damping in the estimate.

The result carries one condition. Classification, unlike location,
depends on how each printed map's hue directions are registered to the
motion-capture axes: a reflection of the sensed vectors changes the sign
of $\det\hat{\mathbf{J}}$ and exchanges saddles with nodes, while leaving
the estimated critical point where it was. The two maps were printed and
placed under different frame conventions (reflections, rotations, and the
camera's mirroring among them). The results above apply to each map the
registration used in the author's own field reconstructions of
Chapter~\ref{ch:tools}, which negates the sensed $v$ component for the
saddle map only. The controller applied no such registration, since it
used the estimate only for location, and its unregistered Jacobian read
the saddle map as a stable node or spiral throughout the terminal window.
A deployment that reports feature type therefore needs the sensor frame
registered to the navigation frame, a requirement that location alone
does not impose.

\section{Attraction Primitive}
\label{sec:ch6:attraction}

\subsection{Definition}
\label{sec:ch6:attraction_def}

% Ported from: cwaight_systems_paper.tex Sec. II-D.1, eq. 10; thesis.tex
% "Critical Point Attraction" (lines 1979-1993).
The three-robot cluster centroid is
$\mathbf{p}_c = \tfrac{1}{3}\sum_{i=1}^{3}\mathbf{p}_i$. To navigate
$\mathbf{p}_c$ onto the estimated critical point, a proportional control
law is applied:
\begin{equation}
\dot{\mathbf{p}}_c = k(\mathbf{p}^* - \mathbf{p}_c), \qquad k>0.
\label{eq:ch6:attraction}
\end{equation}
The mapping from this cluster-level velocity command to individual robot
velocities is handled by the three-robot SAS cluster-space controller of
Chapter~\ref{ch:estimation} and Appendix~\ref{app:kin3}.

\subsection{Stability}
\label{sec:ch6:attraction_stability}

% Ported from: cwaight_systems_paper.tex Sec. II-D.1 (lines 160-164);
% thesis.tex "Critical Point Attraction" (lines 1987-1992).
\begin{proposition}
\label{thm:ch6:attraction_stability}
Under \eqref{eq:ch6:attraction} and a stationary estimate
($\dot{\mathbf{p}}^*=\mathbf{0}$), the centroid converges globally and
exponentially to $\mathbf{p}^*$ with time constant $\tau = 1/k$.
\end{proposition}
\begin{proof}
Define the error $\mathbf{e} = \mathbf{p}_c - \mathbf{p}^*$. Since
$\dot{\mathbf{p}}^*=\mathbf{0}$, $\dot{\mathbf{e}} = \dot{\mathbf{p}}_c =
-k\mathbf{e}$, so $\mathbf{e}(t) = \mathbf{e}(0)e^{-kt}$, independent of
$\mathbf{e}(0)$.
\end{proof}

On an exactly affine field the estimate $\mathbf{p}^*$ is exact at every
measurement, hence automatically stationary in the noise-free case
(Section~\ref{sec:ch6:formation_model}). Sensor noise and formation
conditioning perturb $\mathbf{p}^*$ between control cycles and are
treated separately in Section~\ref{sec:ch6:error_analysis}; actuator
dynamics and command limits, not part of this kinematic model, are also
addressed there.

\subsection{Simulation Verification}
\label{sec:ch6:attraction_sim}

% Ported from: cwaight_systems_paper.tex Sec. IV, IV-A (lines 239-256),
% Table II (lines 259-280); thesis.tex "Convergence to Critical Points"
% (lines 2074-2119).
Simulation trials used the robot dynamics and simulation environment of
Chapter~\ref{ch:tools}, with a cluster of three robots in an equilateral
triangular formation of side length 0.33~m and attraction gain
$k=1$~s$^{-1}$ in \eqref{eq:ch6:attraction}. Six noise-free analytical
fields, one for each canonical critical point type (Chapter~%
\ref{ch:preliminaries}; field formulas in Appendix~\ref{app:fields}),
were used for trials. After verifying performance in noise-free
environments, trials were also run on two fields reconstructed from
robot-sensed measurements collected on the hardware testbed (vortex and
saddle), incorporating realistic sensor noise.

For each of the eight fields, the cluster was initialized at a random
position drawn uniformly from a 1~m~$\times$~1~m box centered on the true
critical point and run for 100 time steps at 10~Hz. This was repeated
1000 times per field; the final centroid position of each trial gave
systematic bias (mean error) and precision (standard deviation). A trial
was counted as convergent if the trajectory remained bounded and
terminated within the stiction-limited neighborhood of the critical
point.

The cluster converged in 100\% of the 1000 trials for every field tested.
On the six analytical fields, systematic bias remained below 0.001~m and
precision ranged from 0.014 to 0.016~m. On the two reconstructed fields,
systematic bias was 0.028~m for the vortex and 0.005~m for the saddle,
with precision near 0.016~m in both cases (Table~\ref{tab:ch6:sim_attraction}).

\begin{figure}[h]
  \centering
  \includegraphics[width=0.75\textwidth]{ch6_figure_5.png}
  \caption{Successful navigation to critical points from random starting
  positions across six canonical vector field types.}
  \label{fig:ch6:attraction_sim}
\end{figure}
% FIGURE SOURCE: Paper_Writing/Systems Paper/cwaight_systems_paper/figure_5.png

\begin{table}[h]
\footnotesize
\centering
\caption{Simulation results across eight field types (six analytical
plus two reconstructed). Every field achieved 100\% convergence over
1000 trials.}
\label{tab:ch6:sim_attraction}
\setlength{\tabcolsep}{4pt}
\begin{tabular}{lccc}
\toprule
\textbf{Field} & \textbf{Avg.\ Rest Point (x, y) m} & \textbf{Bias (m)} & \textbf{$\sigma$ (m)} \\
\midrule
Sink & (0.0006, $-$0.0001) & 0.0006 & 0.0156 \\
Source & ($-$0.0005, $-$0.0001) & 0.0006 & 0.0140 \\
Vortex & ($-$0.0002, $-$0.0000) & 0.0002 & 0.0163 \\
Sinking Vortex & ($-$0.0002, $-$0.0004) & 0.0004 & 0.0145 \\
Spewing Vortex & (0.0001, 0.0006) & 0.0006 & 0.0157 \\
Saddle & (0.0004, $-$0.0001) & 0.0004 & 0.0157 \\
Vortex (reconstructed) & (0.0206, $-$0.0187) & 0.0278 & 0.0162 \\
Saddle (reconstructed) & ($-$0.0048, $-$0.0007) & 0.0048 & 0.0157 \\
\bottomrule
\end{tabular}
\end{table}

\subsection{Hardware Verification}
\label{sec:ch6:attraction_hw}

% Ported from: cwaight_systems_paper.tex Sec. V-B, V-C (lines 352-403);
% thesis.tex "Attraction Controller Results" (lines 2216-2243).
Hardware validation used the Decabot testbed of Chapter~\ref{ch:tools},
focused on two field types spanning the key challenges of vector field
navigation: rotational dynamics (vortex) and hyperbolic instability
(saddle); the saddle is the only canonical field type here that contains
separatrices. For each field, eight starting locations were selected
around the perimeter of the workspace and the attraction controller
\eqref{eq:ch6:attraction} was run 10 times from each location using the
three-robot equilateral triangle formation with a 0.33~m side length,
yielding 80 trials per field. Two saddle trials and one vortex trial were
lost to file management errors, giving 157 usable trials (79 vortex, 78
saddle).

The cluster converged to the critical point in 100\% of the 157 hardware
trials. In the vortex field, systematic bias was 0.012~m with precision
of 0.009~m; in the saddle field, systematic bias was 0.005~m with
precision of 0.014~m (Table~\ref{tab:ch6:hw_attraction}). The narrow
standard deviation bounds across the 10 repetitions per starting position
confirm that the residual scatter reflects sensor noise rather than
uncontrolled experimental variation.

\begin{figure}[h]
  \centering
  \includegraphics[width=0.75\textwidth]{ch6_allrunsadjusted.png}
  \caption{Convergence trajectories from eight starting positions to the
  center of the vortex (a) and saddle (b) using three robots.}
  \label{fig:ch6:attraction_hw}
\end{figure}
% FIGURE SOURCE: Paper_Writing/Systems Paper/cwaight_systems_paper/allrunsadjusted.png

\begin{figure}[h]
  \centering
  \includegraphics[width=\textwidth]{ch6_fig9.png}
  \caption{Distance to center over time, with mean trajectory and
  standard deviation bounds.}
  \label{fig:ch6:attraction_hw_distance}
\end{figure}
% FIGURE SOURCE: Paper_Writing/Systems Paper/cwaight_systems_paper/fig9.png

\begin{table}[h]
\footnotesize
\centering
\caption{Hardware experimental results. Saddle: 78 runs. Vortex: 79 runs.
Bias measures mean offset from the true critical point; precision
($\sigma$) measures repeatability.}
\label{tab:ch6:hw_attraction}
\setlength{\tabcolsep}{4pt}
\begin{tabular}{lccc}
\toprule
\textbf{Field} & \textbf{Avg.\ Rest Point (x, y) m} & \textbf{Bias (m)} & \textbf{$\sigma$ (m)} \\
\midrule
Saddle (printed map) & (0.004, $-$0.001) & 0.005 & 0.014 \\
Vortex (printed map) & ($-$0.004, 0.012) & 0.012 & 0.009 \\
\bottomrule
\end{tabular}
\end{table}

\section{Orbital Primitive}
\label{sec:ch6:orbital}

\subsection{Definition}
\label{sec:ch6:orbital_def}

% Ported from: cwaight_systems_paper.tex Sec. II-D.2, eq. 11; thesis.tex
% "Orbital Navigation" (lines 1994-2014).
To maintain a circular orbit of radius $r_d$ around $\mathbf{p}^*$,
define the radial vector $\mathbf{r} = \mathbf{p}^* - \mathbf{p}_c$ with
magnitude $r = \lVert\mathbf{r}\rVert$. The unit radial vector is
$\hat{\mathbf{r}} = \mathbf{r}/r$, and the tangent vector
$\hat{\mathbf{r}}_\perp$ is obtained by a $-90^\circ$ rotation. The
orbital velocity command is
\begin{equation}
\dot{\mathbf{p}}_c = k_t\hat{\mathbf{r}}_\perp + k_r(r-r_d)\hat{\mathbf{r}},
\label{eq:ch6:orbital}
\end{equation}
where $k_t$ is the commanded orbital (tangential) speed, whose sign fixes
direction, and $k_r>0$ drives radial convergence.

\subsection{Stability}
\label{sec:ch6:orbital_stability}

% Ported from: cwaight_systems_paper.tex Sec. II-D.2 (lines 173-174);
% thesis.tex "Orbital Navigation" (lines 2007-2013).
\begin{proposition}
\label{thm:ch6:orbital_stability}
Under \eqref{eq:ch6:orbital} and a stationary estimate, the radial error
$e_r = r-r_d$ decays exponentially to zero, and as $r\to r_d$ the angular
velocity stabilizes to $\omega = k_t/r_d$, producing steady circular
motion.
\end{proposition}
\begin{proof}
Equation~\eqref{eq:ch6:orbital} decouples radial and angular dynamics:
the radial error $e_r=r-r_d$ obeys $\dot e_r = -k_r e_r$, so
$e_r(t)=e_r(0)e^{-k_r t}$, ensuring exponential convergence to $r_d$. As
$r\to r_d$, the angular velocity stabilizes to $\omega=k_t/r_d$,
producing steady circular motion.
\end{proof}

\subsection{Simulation Verification}
\label{sec:ch6:orbital_sim}

% Ported from: cwaight_systems_paper.tex Sec. IV-B (lines 283-311),
% Table III; thesis.tex "Orbital Control" (lines 2120-2163).
Using the simulation setup of Section~\ref{sec:ch6:attraction_sim},
orbital control was evaluated on the same eight fields at six commanded
radii with gains $k_t=0.2$~m/s and $k_r=0.5$~s$^{-1}$ in
\eqref{eq:ch6:orbital}. For each radius, the cluster was positioned at
the desired radius to initialize the run, then the orbital controller
was used for 600 time steps; mean radial error and standard deviation
were computed after discarding the first 100 time steps (10~s) as a
settling transient.

On the six noise-free fields, all six produced identical performance at
each radius, with standard deviations below $2\times10^{-5}$~m
(Table~\ref{tab:ch6:sim_orbital}). This is expected: on an exactly
affine noise-free field the three-robot fit recovers $\mathbf{p}^*$
exactly at every step, so the controller input, and hence the
trajectory, is independent of the field type; the residual standard
deviation is the steady-state radius ripple rather than field variation,
hence the single noise-free column in the table. Mean radial errors had
a slight positive bias that decreased with increasing radius, from
0.146~m at a commanded radius of 0.01~m to 0.042~m at 0.50~m. This
pattern held on the reconstructed vortex, with mean errors from 0.150~m
to 0.027~m, and on the reconstructed saddle, from 0.146~m to
$-0.003$~m. In both reconstructed fields, error standard deviations
increased with commanded radius, from 0.012~m to 0.061~m.

No figure in the source paper covers this simulated-orbital experiment;
it is reported there only as Table~\ref{tab:ch6:sim_orbital}.


\begin{table}[h]
\centering
\caption{Simulated orbital control across eight field types. Negative
error indicates the cluster orbited inside the commanded radius.}
\label{tab:ch6:sim_orbital}
\footnotesize
\setlength{\tabcolsep}{3pt}
\begin{tabular}{cccc}
\toprule
\textbf{Radius (m)} & \textbf{Noise-free (6 fields)} & \textbf{Vortex (reconstructed)} & \textbf{Saddle (reconstructed)} \\
 & \textbf{Error (m) $\pm$ std} & \textbf{Error (m) $\pm$ std} & \textbf{Error (m) $\pm$ std} \\
\midrule
0.01 & 0.146 $\pm$ $3{\times}10^{-6}$ & 0.150 $\pm$ 0.015 & 0.146 $\pm$ 0.012 \\
0.10 & 0.109 $\pm$ $7{\times}10^{-6}$ & 0.109 $\pm$ 0.013 & 0.098 $\pm$ 0.015 \\
0.20 & 0.081 $\pm$ $1{\times}10^{-5}$ & 0.104 $\pm$ 0.014 & 0.072 $\pm$ 0.024 \\
0.30 & 0.063 $\pm$ $1{\times}10^{-5}$ & 0.094 $\pm$ 0.018 & 0.059 $\pm$ 0.042 \\
0.40 & 0.050 $\pm$ $2{\times}10^{-5}$ & 0.067 $\pm$ 0.026 & 0.035 $\pm$ 0.054 \\
0.50 & 0.042 $\pm$ $2{\times}10^{-5}$ & 0.027 $\pm$ 0.038 & $-$0.003 $\pm$ 0.061 \\
\bottomrule
\end{tabular}
\end{table}

\subsection{Hardware Verification}
\label{sec:ch6:orbital_hw}

% Ported from: cwaight_systems_paper.tex Sec. V-B, V-D (lines 359-361,
% 407-418), Table V; thesis.tex "Orbital Control Results"
% (lines 2244-2280).
Using the same testbed and field pair as
Section~\ref{sec:ch6:attraction_hw}, the orbiting controller
\eqref{eq:ch6:orbital} was tested at six commanded radii matching the
simulation distances, each run for 3000 cycles at 10~Hz, more than three
complete orbits. Mean radial error and standard deviation were recorded
over the full run, including the initial transient (unlike the
simulation statistics of Table~\ref{tab:ch6:sim_orbital}, which discard a
10~s settling transient).

In the vortex field, mean radial errors decreased from 0.246~m at a
commanded radius of 0.01~m to $-0.035$~m at 0.50~m, with standard
deviations ranging from 0.012 to 0.053~m. In the saddle field, mean
radial errors ranged from 0.114~m to 0.060~m over the same commanded
radii, with standard deviations increasing from 0.011 to 0.095~m
(Table~\ref{tab:ch6:hw_orbital}).

\begin{figure}[h]
  \centering
  \includegraphics[width=\textwidth]{ch6_orbit040_analytical_vs_real.png}
  \caption{Orbital trajectory comparison at a commanded radius of
  0.40~m: analytical prediction vs. hardware performance in the saddle
  field.}
  \label{fig:ch6:orbit_hw}
\end{figure}
% FIGURE SOURCE: Paper_Writing/Systems Paper/cwaight_systems_paper/orbit040_analytical_vs_real.png

\begin{table}[h]
\centering
\caption{Hardware orbital control results. Negative error indicates the
cluster orbited inside the commanded radius.}
\label{tab:ch6:hw_orbital}
\begin{tabular}{ccc}
\toprule
\textbf{Radius (m)} & \textbf{Vortex (printed map)} & \textbf{Saddle (printed map)} \\
 & \textbf{Error (m) $\pm$ std} & \textbf{Error (m) $\pm$ std} \\
\midrule
0.01 & 0.246 $\pm$ 0.053 & 0.114 $\pm$ 0.011 \\
0.10 & 0.190 $\pm$ 0.012 & 0.072 $\pm$ 0.018 \\
0.20 & 0.143 $\pm$ 0.019 & 0.058 $\pm$ 0.037 \\
0.30 & 0.087 $\pm$ 0.015 & 0.089 $\pm$ 0.059 \\
0.40 & 0.032 $\pm$ 0.020 & 0.090 $\pm$ 0.079 \\
0.50 & $-$0.035 $\pm$ 0.023 & 0.060 $\pm$ 0.095 \\
\bottomrule
\end{tabular}
\end{table}

\section{Comparison with Zeroth-Order Primitives}
\label{sec:ch6:comparison}

% Ported from: cwaight_systems_paper.tex Sec. VI-B "Comparison with
% Other Primitives" (lines 462-476); thesis.tex "Comparison with Other
% Primitives" (lines 2328-2345), reframed against Chapter 5.
The modular control architecture of Chapter~\ref{ch:estimation} allows
navigation controllers to be interchanged while the cluster-space and
robot controller layers remain fixed, so performance differences reflect
algorithmic rather than implementation variation. This isolates the
contribution of location estimation: the orbital primitive of
Section~\ref{sec:ch6:orbital} is compared against the vector-sum and
vector-to-scalar primitives of Chapter~\ref{ch:zeroth} in the vortex
field. The vector-sum primitive commands motion in the direction of the
average perceived flow; the vector-to-scalar primitive \cite{38},
modified here to move tangentially to the direction of steepest descent,
reduces the flow to its magnitude.

Both zeroth-order alternatives exhibit outward spiral drift in the
vortex field (Figure~\ref{fig:ch6:primitive_comparison}). Both derive
their velocity commands from direction-only guidance; neither has a
mechanism to detect or correct radial error with respect to a specific
point, so centrifugal drift goes uncorrected. The vector-sum drift has
been confirmed in hardware \cite{39}. In the saddle field the
comparison is more decisive: the field vectors do not all point toward
the critical point, so the vector-sum controller cannot produce an
orbital trajectory at all, and the vector-to-scalar technique would
still drift. The orbital primitive of this chapter avoids both failure
modes because it holds an explicit estimate of $\mathbf{p}^*$ and closes
the loop on the radial error $r-r_d$ directly, rather than acting on
direction alone.

\begin{figure}[h]
  \centering
  \includegraphics[width=\textwidth]{ch6_orbit_demonstration.png}
  \caption{Orbital trajectories: vector-sum primitive versus the
  proposed critical point controller.}
  \label{fig:ch6:primitive_comparison}
\end{figure}
% FIGURE SOURCE: Paper_Writing/Systems Paper/cwaight_systems_paper/orbit_demonstration.png

\section{Error Analysis}
\label{sec:ch6:error_analysis}

% Ported from: thesis.tex "Error Analysis" (lines 2283-2327; primary
% source for the robot-dynamics, equilateral-triangle kappa(A)=7.42, and
% Jacobian-conditioning material); cwaight_systems_paper.tex Discussion
% Sec. VI-A (sigma_p ~ ||J^{-1}|| sigma_v / sqrt(3) noise-propagation
% paragraph, line 454, not carried into thesis.tex's condensed version).
Estimation and control errors arise from four sources: robot dynamics,
sensor noise, cluster geometry, and field conditioning.

Perfect convergence does not occur in simulation. However, repeating
experiments with no stiction, no maximum velocity, and
$\alpha_{\text{mom}}=0$ (instantaneous velocity response,
Chapter~\ref{ch:tools}) confirms convergence to the controller's
$10^{-6}$~m command deadband, with the critical point estimate itself
exact to machine precision. This confirms the residual variance is due
to robot dynamics, not the estimation method.

Positive radial bias at small orbital radii is also a result of robot
dynamics. At small radii, the commanded angular rate $k_t/r_d$ is high
relative to the radial correction bandwidth $k_r$, and the centripetal
demand ($k_t^2/r$) pushes the cluster outward faster than the radial
controller can correct; as the commanded radius increases, this effect
diminishes. More aggressive radial gains would reduce this error, at the
cost of increased sensitivity to noise in the critical point estimate.

Comparing the noise-free and reconstructed fields in simulation isolates
the effect of noisy sensor data: sensor noise is the driving factor of
the simulation-reality gap, with sensing errors propagating directly
through the critical point estimation equations
\eqref{eq:ch6:jacobian_offset}--\eqref{eq:ch6:pstar}. To first order, the
centered equilateral fit maps per-component sensor noise $\sigma_v$ into
estimate scatter
\begin{equation}
\sigma_p \approx \frac{\lVert\mathbf{J}^{-1}\rVert\,\sigma_v}{\sqrt{3}};
\label{eq:ch6:sigma_p}
\end{equation}
the calibration errors of Chapter~\ref{ch:tools} give
$\sigma_v \approx 0.02$~m/s, predicting roughly 0.03~m of scatter on the
printed vortex. This matches the observed 0.009--0.016~m precision after
closed-loop time-averaging, while the 0.012--0.028~m bias reflects
systematic distortion that averaging cannot remove.

To avoid errors from robots becoming collinear (the condition number of
$\mathbf{A}$ approaching infinity), an equilateral triangle was chosen
as the cluster shape: among 200{,}000 random formations of the same size
centered on the sampled region, none beat the equilateral triangle's
$\kappa(\mathbf{A})=7.42$. This optimality is not translation-invariant
(the same formation centered elsewhere has a different condition
number), so it is stated for a formation centered on the sampled region,
as here.

Further errors arise from poor conditioning of the estimated field
Jacobian $\kappa(\mathbf{J})$, directly affecting the critical point
estimate \eqref{eq:ch6:pstar}. The condition number grows when the
Jacobian eigenvalues have very different magnitudes, as occurs near
highly asymmetric nodes or saddle points with disparate stable and
unstable manifold strengths, and diverges as $\det(\mathbf{J})\to0$,
corresponding to non-hyperbolic critical points or near-uniform flow.
The negative bias at the largest commanded radii traces to degradation
of the reconstructed field near the edge of the mapped region: there the
estimated critical point distance is systematically inflated, so the
radial controller pulls the cluster inside the commanded radius, the
same range-dependent error amplification captured by $\kappa(\mathbf{J})$,
since far from the critical point the sampled patch approaches uniform
flow and small field errors produce large estimate shifts. Unlike
$\kappa(\mathbf{A})$, this factor is a property of the field, not the
formation, and is not under the designer's control.

Across all four error sources, estimation errors shift the equilibrium
location of both controllers but preserve the stability established in
Sections~\ref{sec:ch6:attraction_stability} and~\ref{sec:ch6:orbital_stability}.

\section{Summary}
\label{sec:ch6:summary}

% Ported from: thesis.tex "Conclusion" (lines 2388-2397), trimmed to
% this chapter's own results; forward pointers to Chapter 7 and to
% future-work material move to Chapters 9 and 10.
This chapter presented a distributed estimation framework that uses
instantaneous measurements from three robots to recover both the
location and the type of a nearby critical point, and two control
primitives, attraction and orbit, built on that estimate. Simulation
across six canonical field types confirmed convergence of both
primitives (Sections~\ref{sec:ch6:attraction_sim},
\ref{sec:ch6:orbital_sim}). In 157 hardware attraction trials, the
system achieved a 100\% success rate with position errors of 0.012~m or
less (Section~\ref{sec:ch6:attraction_hw}). Orbital control experiments
maintained commanded radii within 0.25~m error at small radii and 0.1~m
at larger radii (Section~\ref{sec:ch6:orbital_hw}). Against the
zeroth-order primitives of Chapter~\ref{ch:zeroth}, the explicit
critical point estimate removes the outward drift both alternatives
exhibit on a rotational feature (Section~\ref{sec:ch6:comparison}).
